% Aalok Marathi journal structure starter.
\documentclass[12pt,a4paper]{article}
\usepackage{fontspec}
\IfFontExistsTF{Noto Serif Devanagari}{\setmainfont{Noto Serif Devanagari}}{\setmainfont{FreeSerif}}
\usepackage[margin=2.2cm]{geometry}
\usepackage{booktabs,tabularx,microtype,xcolor,hyperref}
\hypersetup{colorlinks=true,urlcolor=blue!50!black,linkcolor=blue!50!black}
\setlength{\parindent}{0pt}
\setlength{\parskip}{0.55em}
\newcommand{\marathi}[1]{#1}
\title{A reproducible workflow for a Marathi journal article}
\author{Anita Deshmukh\\\small Example Institute of Language Technology\\\small\texttt{anita@example.org}}
\date{Volume 12, Issue 2 \textbullet{} 2026}
\begin{document}
\maketitle

\begin{abstract}
  This neutral starter demonstrates the compact article structure used by a
  Marathi-language journal. The example records a reproducible editorial
  workflow and keeps every value editable.
\end{abstract}
\textbf{Keywords:} Marathi publishing; reproducibility; editorial workflow

\section{Introduction}
Small journals benefit from a clear separation between the author's content
and the production metadata added by an editorial office. A structured source
also makes it easier to preserve the language, citations, and tables when an
article moves between contributors.

\section{लेखाची रचना}
हा परिच्छेद मराठी मजकूर, शीर्षके आणि देवनागरी अंक संपादनासाठी जागा दाखवतो.
लेखकाने येथे स्वतःचा अभ्यास, पद्धत आणि निष्कर्ष लिहावा.

\section{Method}
The workflow has three stages: define the article metadata, write the
argument in numbered sections, and run a final language and reference check.
Table~\ref{tab:stages} shows a compact audit record.

\begin{table}[ht]
  \centering
  \caption{Example editorial audit record.}
  \label{tab:stages}
  \begin{tabularx}{0.9\linewidth}{@{}lXr@{}}
    \toprule
    Stage & Evidence captured & Checks \\
    \midrule
    Metadata & title, author, affiliation, dates & 4 \\
    Manuscript & sections, figures, tables & 7 \\
    Release & PDF, source, accessibility review & 5 \\
    \bottomrule
  \end{tabularx}
\end{table}

\section{Conclusion}
The starter is intentionally neutral: replace the example text, add the
journal's current front matter, and verify the official submission rules
before publication.

\end{document}
